\documentclass[12pt,a4paper,reqno]{amsart}

% language
\usepackage[greek,english]{babel}
\usepackage[utf8]{inputenc}
\usepackage{alphabeta}

% change default names to greek
\addto\captionsenglish{
    \renewcommand{\contentsname}{Περιεχόμενα}
    \renewcommand{\refname}{Βιβλιογραφία}
    \renewcommand{\datename}{Ημερομηνία:}
    \renewcommand{\urladdrname}{Ιστοσελίδα}
}

% math 
\usepackage{amsmath,amsthm,amssymb,amscd}

% font
\usepackage[cal=euler]{mathalfa}
\usepackage{libertinus-type1}
% \usepackage{txfonts} % for upright greek letters
\usepackage{bm} % for bold symbols
\usepackage{bbm} % for the simply-looking bb symbols

% miscellaneous 
\usepackage{changepage} %for indenting environments
\usepackage{csquotes} % example: \textcquote{}
\usepackage{blkarray}
\setcounter{MaxMatrixCols}{20} % default for pmatrix is 10!!
\usepackage{ytableau}
\usepackage{array} %needed to increase the vertical length in a tabular

% drawing
\usepackage{tikz,tikz-cd}
\usetikzlibrary{shapes.misc, patterns, matrix, calc, intersections,positioning,arrows.meta}
\usepackage{graphics,graphicx}
\usepackage{float} % provides enhanced control and customization options for floating objects such as figures and tables

% colors
\usepackage{xcolor}
\definecolor{darkcandyapplered}{rgb}{0.64, 0.0, 0.0}
\definecolor{midnightblue}{rgb}{0.1, 0.1, 0.44}
\definecolor{mylightblue}{HTML}{336699}
\definecolor{burntorange}{rgb}{0.8, 0.33, 0.0}
\definecolor{iceberg}{rgb}{0.44, 0.65, 0.82}
\definecolor{applegreen}{rgb}{0.55, 0.71, 0.0}
\definecolor{canaryyellow}{rgb}{1.0, 0.94, 0.0}
\definecolor{lightgray}{rgb}{0.83, 0.83, 0.83}

% hrefs
\usepackage{hyperref}
\usepackage[noabbrev,capitalize]{cleveref}
\hypersetup{
    pdftoolbar=true,        
    pdfmenubar=true,        
    pdffitwindow=false,     
    pdfstartview={FitH},  % fits the width of the page to the window
    pdftitle={},
    pdfauthor={},
    pdfsubject={},
    pdfkeywords={},
    pdfnewwindow=true,  % links in new window
    colorlinks=true,  % false: boxed links; true: colored links
    linkcolor=darkcandyapplered,   % color of internal links
    citecolor=midnightblue,  % color of links to bibliography
    urlcolor=cyan,  % color of external links
    linktocpage=true  % changes the links from the section body to the page number
    }

% geometry
\textwidth=16cm 
\textheight=21cm 
\hoffset=-55pt 
\footskip=25pt

% thm envs (you might need to change the path)
\theoremstyle{definition}
\newtheorem{exercise}{Άσκηση}
\newtheorem{solution}{Λύση}
\newtheorem*{remark}{Παρατήρηση}
\newtheorem*{example}{Παράδειγμα}

% math ops 
% fields
\newcommand{\NN}{\mathbbmss{N}} 
\newcommand{\ZZ}{\mathbbmss{Z}} 
\newcommand{\QQ}{\mathbbmss{Q}} 
\newcommand{\RR}{\mathbbmss{R}} 
\newcommand{\CC}{\mathbbmss{C}} 
\newcommand{\KK}{\mathbbmss{K}} 
\newcommand{\FF}{\mathbbmss{F}} 
\newcommand{\HH}{\mathbbmss{H}} 

% symmetric group
\newcommand{\fS}{\mathfrak{S}}  
\newcommand{\fp}{\mathfrak{p}}  
\newcommand{\fm}{\mathfrak{m}}  

% calligraphic 
\newcommand{\aA}{\mathcal{A}} 
\newcommand{\bB}{\mathcal{B}}
\newcommand{\cC}{\mathcal{C}}
\newcommand{\dD}{\mathcal{D}}
\newcommand{\eE}{\mathcal{E}}
\newcommand{\fF}{\mathcal{F}}
\newcommand{\hH}{\mathcal{H}}
\newcommand{\iI}{\mathcal{I}}
\newcommand{\lL}{\mathcal{L}}
\newcommand{\nN}{\mathcal{N}}
\newcommand{\oO}{\mathcal{O}}
\newcommand{\pP}{\mathcal{P}}
\newcommand{\sS}{\mathcal{S}}
\newcommand{\mM}{\mathcal{M}}
\newcommand{\uU}{\mathcal{U}}

% bold
\newcommand{\bfa}{\mathbf{a}}
\newcommand{\bfe}{\mathbf{e}}
\newcommand{\bfF}{\pmb{F}}
\newcommand{\bfR}{\pmb{R}}
\newcommand{\bfv}{\mathbf{v}}
%\newcommand{\bfx}{\bm{x}}
%\newcommand{\bfx}{\mathbf{x}} 
\newcommand{\bfx}{\pmb{x}}
\newcommand{\bfX}{\pmb{X}}
\newcommand{\bfy}{\pmb{y}}
\newcommand{\bfz}{\pmb{z}}

% roman
\newcommand{\rmA}{\mathrm{A}}
\newcommand{\rmB}{\mathrm{B}}
\newcommand{\rmC}{\mathrm{C}}
\newcommand{\rmD}{\mathrm{D}} 
\newcommand{\rmH}{\mathrm{H}} 
\newcommand{\rmh}{\mathrm{h}} 
\newcommand{\rmI}{\mathrm{I}} 
\newcommand{\rmJ}{\mathrm{J}} 
\newcommand{\rmK}{\mathrm{K}}
\newcommand{\rmM}{\mathrm{M}}
\newcommand{\rmP}{\mathrm{P}}  
\newcommand{\rmp}{\mathrm{p}}  
\newcommand{\rmQ}{\mathrm{Q}}  
\newcommand{\rmR}{\mathrm{R}}
\newcommand{\rmS}{\mathrm{S}}
\newcommand{\rmT}{\mathrm{T}}
\newcommand{\rmU}{\mathrm{U}}
\newcommand{\rmV}{\mathrm{V}}
\newcommand{\rmY}{\mathrm{Y}}
\newcommand{\rmZ}{\mathrm{Z}}
\newcommand{\rmz}{\mathrm{z}}

% arrows and symbols 
\renewcommand{\to}{\rightarrow}
\newcommand{\toto}{\longrightarrow}
\newcommand{\mapstoto}{\longmapsto}
\newcommand{\then}{\Rightarrow}
\newcommand{\IFF}{\Leftrightarrow}
\newcommand{\tl}{\tilde}
\newcommand{\wtl}{\widetilde}
\newcommand{\ol}{\overline}
\newcommand{\ul}{\underline}
\newcommand{\oldemptyset}{\emptyset}
\renewcommand{\emptyset}{\varnothing}
\DeclareMathSymbol{\Arg}{\mathbin}{AMSa}{"39} % for arguments 
\newcommand{\onto}{\ensuremath{\twoheadrightarrow}}
\newcommand{\tle}{\trianglelefteq}
\newcommand{\tge}{\trianglerighteq}

% custom ops
\newcommand{\rmKer}{\mathrm{Ker}}
\newcommand{\rmIm}{\mathrm{Im}}
\newcommand{\lcm}{\mathrm{lcm}}
\newcommand{\GL}{\mathrm{GL}}
\newcommand{\ev}{\mathrm{ev}}
\newcommand{\End}{\mathrm{End}}
\newcommand{\rmid}{\mathrm{id}}
\newcommand{\rmspan}{\mathrm{span}}
\newcommand{\Spec}{\mathrm{Spec}}
\newcommand{\mSpec}{\mathrm{mSpec}}
\newcommand{\rad}{\mathrm{rad}}
\newcommand{\Rad}{\mathrm{Rad}}
\newcommand{\Jac}{\mathrm{Jac}}
\newcommand{\tor}{\mathrm{tor}}
\newcommand{\Ann}{\mathrm{Ann}}
\newcommand{\Hom}{\mathrm{Hom}}

% absolute value symbol
\usepackage{mathtools} 
\DeclarePairedDelimiter\abs{\lvert}{\rvert}%
\DeclarePairedDelimiter\norm{\lVert}{\rVert}%
\makeatletter
\let\oldabs\abs
\def\abs{\@ifstar{\oldabs}{\oldabs*}}

% tensor symbol
\newcommand{\tensor}[1]{%
  \mathbin{\mathop{\otimes}\limits_{#1}}%
}

% permutation cycle notation
\ExplSyntaxOn
\NewDocumentCommand{\cycle}{ O{\;} m }
 {
  (
  \alec_cycle:nn { #1 } { #2 }
  )
 }

\seq_new:N \l_alec_cycle_seq
\cs_new_protected:Npn \alec_cycle:nn #1 #2
 {
  \seq_set_split:Nnn \l_alec_cycle_seq { , } { #2 }
  \seq_use:Nn \l_alec_cycle_seq { #1 }
 }
\ExplSyntaxOff

% setminus symbol
\newcommand{\mysetminusD}{\hbox{\tikz{\draw[line width=0.6pt,line cap=round] (3pt,0) -- (0,6pt);}}}
\newcommand{\mysetminusT}{\mysetminusD}
\newcommand{\mysetminusS}{\hbox{\tikz{\draw[line width=0.45pt,line cap=round] (2pt,0) -- (0,4pt);}}}
\newcommand{\mysetminusSS}{\hbox{\tikz{\draw[line width=0.4pt,line cap=round] (1.5pt,0) -- (0,3pt);}}}
\newcommand{\sm}{\mathbin{\mathchoice{\mysetminusD}{\mysetminusT}{\mysetminusS}{\mysetminusSS}}}

% miscellaneous commands
\newcommand{\defn}[1]{{\color{mylightblue}{#1}}}
\newcommand{\toDo}{{\bf\color{red} TODO}}
\newcommand{\toCite}{{\bf\color{green} CITE}}
\newcommand*{\vertbar}{\rule[-1ex]{0.5pt}{2.5ex}} % for matrices with column vectors
\newcommand*{\horzbar}{\rule[.5ex]{2.5ex}{0.5pt}} % for matrices with row vectors
\newcommand{\myblue}[1]{{\color{iceberg}{#1}}}
\newcommand{\myorange}[1]{{\color{burntorange}{#1}}}
\newcommand{\mygreen}[1]{{\color{applegreen}{#1}}}
\newcommand{\myred}[1]{{\color{darkcandyapplered}{#1}}}

% ferrer's diagram
\newcommand{\fdiagram}[1]{
    \begin{tikzpicture}[scale=.7]
        \fill foreach \Z [count=\Y] in {#1}
        {foreach \X in {1,...,\Z} 
        {(\X,-\Y) circle[radius=3pt]}};
    \end{tikzpicture}
}

%
\usepackage{pgfplots}
\pgfplotsset{compat=1.18}

%
\makeatletter
\def\mathcolor#1#{\@mathcolor{#1}}
\def\@mathcolor#1#2#3{%
  \protect\leavevmode
  \begingroup
    \color#1{#2}#3%
  \endgroup
}
\makeatother

% 
\newenvironment{nouppercase}{%
  \let\uppercase\relax%
  \renewcommand{\uppercasenonmath}[1]{}}{}

% titlepage
\title{226: Θεωρία Δακτυλίων και Modules \\ Ασκήσεις με Ενδεικτικές Λύσεις}
\author[Β.~Δ. Μουστακας]{Βασίλης Διονύσης Μουστάκας \\ Πανεπιστήμιο Κρήτης}
\date{17 Μαρτίου 2026}
% \urladdr{\href{https://sites.google.com/view/vasmous}{https://sites.google.com/view/vasmous}}

\begin{document}

\begingroup
\def\uppercasenonmath#1{} % this disables uppercase title
\let\MakeUppercase\relax % this disables uppercase authors
\maketitle
\endgroup

\setcounter{exercise}{15}
% \setcounter{theorem}{5}
\setcounter{equation}{3}
\renewcommand{\thesubsection}{\thesection.\arabic{subsection}}

Σε ότι ακολουθεί υποθέτουμε ότι $R$ είναι ένας δακτύλιος και $n$ είναι ένας θετικός ακέραιος, εκτός αν αναφέρεται διαφορετικά.

\begin{exercise}{\rm(\cite[Άσκηση~10.1.8]{DF04})}
  \label{ex:DF_10.1.8}
  Έστω $Μ$ ένα $R$-πρότυπο. Το $m \in M$ ονομάζεται \defn{στοιχείο στρέψης} (torsion element) αν 
  \[
  rm = 0_m
  \]
  για κάποιο μη μηδενικό $r \in R$. Να δείξετε τα εξής.
  \begin{itemize}
    \item[(α)] Αν ο $R$ είναι ακέραια περιοχή, τότε το σύνολο $M_{\tor}$ όλων των στοιχείων στρέψης του $M$ είναι υποπρότυπο του $M$.
    \item[(β)] Το (α) παύει να ισχύει αν ο $R$ δεν είναι ακέραια περιοχή.
    \item[(γ)] Αν ο $R$ έχει μηδενοδιαιρέτες και $M \neq 0$, τότε $M_\tor \neq 0$.
  \end{itemize}
\end{exercise}

\begin{proof}[Λύση]
  \leavevmode
  \begin{itemize}
    \item[(α)] Υποθέτουμε ότι ο $R$ είναι ακέραια περιοχή και θα επαληθεύσουμε τον Ορισμό 4.2. Το $M_\tor$ περιέχει το μηδέν, διότι $1_R0_M = 0_M$. Το $M_\tor$ είναι κλειστό ως προς την πρόσθεση. Πράγματι, έστω $a, b \in M_\tor$. Τότε $ra = sb = 0_M$, για κάποια μη μηδενικά $r, s \in R$. Επειδή ο $R$ είναι ακέραια περιοχή έπεται ότι $rs \neq 0_R$ και 
    \[
    rs(a+b) = rsa + rsb = s(\underbrace{ra}_{= \, 0_M}) + r(\underbrace{sb}_{= \, 0_M}) = 0_M.
    \]
    Επομένως, $a + b \in M_\tor$. 

    Τέλος, το $M_\tor$ είναι αναλλοίωτο ως προς τη δράση του $R$. Πράγματι, έστω $a \in M_\tor$ και $r \in R$. Τότε $sa = 0_M$, για κάποιο μη μηδενικό $s \in R$. Συνεπώς, 
    \[
    s(ra) = r(\underbrace{sa}_{= \, 0_M}) = 0_M
    \]
    και γι αυτό $ra \in M_\tor$.
    \item[(β)] Θεωρούμε το κανονικό πρότυπο $\ZZ_6$. Το σύνολο των στοιχείων στρέψης του δεν είναι κλειστό ως προς την πρόσθεση. Πράγματι, έχουμε $(\ZZ_6)_\tor = \{0, 2, 3, 4 \}$ (γιατί;), ενώ
    \[
    2 + 3 = 5 \notin (\ZZ_6)_\tor.
    \]
    \item[(γ)] Θεωρούμε ένα μη μηδενικό $m \in M$ και έναν μηδενοδιαιρέτη $r \in R$. Τότε $r \neq 0$ και $rs = 0_R$, για κάποιο μη μηδενικό $s \in R$. Αν $rm = 0_M$, τότε $m \in M_\tor$ και γι αυτό $M_\tor \neq 0$. Διαφορετικά, 
    \[
    s(rm) = (\underbrace{rs}_{= \, 0_R})m = 0_Rm = 0_M,
    \]
    που σημαίνει ότι $rm \in M_\tor$ και γι αυτό $M_\tor \neq 0$.
  \end{itemize}
\end{proof}

\begin{exercise}{\rm(\cite[Ασκήσεις~10.1.9~και~10.1.11]{DF04})}
  \label{ex:DF_10.1.9-10.1.11}
  Έστω $M$ ένα $R$-πρότυπο. Το σύνολο 
  \[
  \Ann_R(M) \coloneqq \{r \in R : rm = 0_M, \ \text{για κάθε $m \in M$}\}
  \]
  ονομάζεται \defn{μηδενιστής} του $M$ στο $R$. Να δείξετε τα εξής.
  \begin{itemize}
    \item[(α)] Ο μηδενιστής του $R$ στο $M$ είναι ιδεώδες του $R$.
    \item[(β)] Να υπολογιστεί ο μηδενιστής του $\ZZ_{24}\times\ZZ_{15}\times\ZZ_{50}$ στο $\ZZ$.
    \item[(γ)] Αν $I$ είναι ένα ιδεώδες του $R$, τέτοιο ώστε $I \subseteq \Ann_R(M)$, τότε η δράση του $R/I$ στο $M$ που ορίζεται θέτοντας 
    \[
    (r + I)\cdot{m} = rm
    \]
    δίνει στο $M$ τη δομή $R/I$-προτύπου.
  \end{itemize}
\end{exercise}

\begin{proof}[Λύση]
  \leavevmode
  \begin{itemize}
    \item[(α)] Αρκεί να παρατηρήσουμε ότι το $\Ann_R(M)$ είναι ο πυρήνας του ομομορφισμού δακτυλίων 
    \begin{align*}
      R &\to \End_\ZZ(M) \\
      r &\mapsto ( m \mapsto rm).
    \end{align*}
    \item[(β)] Αρχικά, παρατηρούμε ότι 
    \[
    \Ann_{\ZZ}(\ZZ_n) = n\ZZ
    \]
    (γιατί;). Έπειτα, έχουμε $x \in \Ann_\ZZ(\ZZ_{24}\times\ZZ_{15}\times\ZZ_{50})$ αν και μόνο αν
    \[
       (0 + 24\ZZ, 0 + 15\ZZ, 0 + 50\ZZ) = (xa + 24\ZZ, xb + 15\ZZ, xc + 50\ZZ)
    \]
    για κάθε $a, b, c \in \ZZ$, ή ισοδύναμα αν και μόνο αν $x \in \Ann_\ZZ(\ZZ_{24})\cap\Ann_\ZZ(\ZZ_{15})\cap\Ann_\ZZ(\ZZ_{50})$. Άρα, 
    \begin{align*}
    \Ann_\ZZ(\ZZ_{24}\times\ZZ_{15}\times\ZZ_{50}) &=  \Ann_\ZZ(\ZZ_{24})\cap\Ann_\ZZ(\ZZ_{15})\cap\Ann_\ZZ(\ZZ_{50}) \\
    &= 24\ZZ \cap 15\ZZ \cap 50\ZZ \\
    &= \lcm(24, 15, 50)\ZZ \\ 
    &= 600\ZZ.
    \end{align*}
    \item[(γ)] Θα δείξουμε ότι η δράση είναι καλώς ορισμένη. Η επαλήθευση του Ορισμού 4.1 αφήνεται ως άσκηση. Πράγματι, έστω $r + I = s + I$ για κάποια $r, s \in R$. Τότε $r-s \in I \subseteq \Ann_R(M)$ και γι αυτό $(r-s)m = 0_M$, για κάθε $m \in M$, ή ισοδύναμα $rm = sm$, για κάθε $m \in M$ και το ζητούμενο έπεται.
  \end{itemize}
\end{proof}

\begin{exercise}{\rm(\cite[Ασκήσεις~10.3.4~και~10.3.5]{DF04})}
  \label{ex:DF_10.3.4-10.3.5}
  Έστω $M$ ένα $R$-πρότυπο. Το $M$ ονομάζεται 
  \begin{itemize}
    \item \defn{πρότυπο στρέψης} (torsion module) αν $M=M_\tor$.
    \item \defn{ελεύθερο στρέψης} (torsion-free) αν $M_\tor = 0$.
  \end{itemize}
  \begin{itemize}
    \item[(α)] Να δείξετε ότι κάθε πεπερασμένη αβελιανή ομάδα είναι $\ZZ$-πρότυπο στρέψης και να δώσετε παράδειγμα άπειρης αβελιανής ομάδας η οποία είναι $\ZZ$-πρότυπο στρέψης.
    \item[(β)] Υποθέτουμε ότι ο $R$ είναι ακέραια περιοχή. Αν το $M$ είναι πεπερασμένα παραγόμενο πρότυπο στρέψης, τότε να δείξετε ότι $\Ann_R(M) \neq 0$ και να δώσετε παράδειγμα προτύπου στρέψης με τετριμμένο μηδενιστή.
    \item[(γ)] Υποθέτουμε ότι ο $R$ είναι ακέραια περιοχή. Να δείξετε ότι κάθε ελεύθερο $R$-πρότυπο είναι ελεύθερο στρέψης.
  \end{itemize}
\end{exercise}

\begin{proof}[Λύση]
  \leavevmode
  \begin{itemize}
    \item[(α)] Αν $R=\ZZ$, αρκεί να παρατηρήσουμε ότι τα στοιχεία στρέψης του $M$ είναι ακριβώς τα στοιχεία πεπερασμένης τάξης (γιατί;). Για το παράδειγμα, θεωρούμε το $\ZZ$-πρότυπο 
    \[
    M = \bigoplus_{n \ge 1} \ZZ_n.
    \]
    Τότε το $M$ είναι προφανώς άπειρο και αποτελείται από στοιχεία της μορφής $(a_1, a_2,$ $\dots)$, για $a_i = 0 + \ZZ_i$ εκτός από ένα το πολύ πεπερασμένο πλήθος. Αν τα μη μηδενικά στοιχεία βρίσκονται στις θέσεις $i_1, i_2, \dots, i_k$, τότε 
    \[
    \lcm(\abs{a_{i_1}}, \abs{a_{i_2}}, \dots, \abs{a_{i_k}})(a_1, a_2, \dots) = 0_M
    \]
    όπου $\abs{a_i}$ είναι η τάξη του $a_i$ στο $\ZZ_i$ για κάθε $i \ge 1$ (γιατί;). Κατά συνέπεια, το $M$ είναι πρότυπο στρέψης.

    \item[(β)] Υποθέτουμε ότι ο $R$ είναι ακέραια περιοχή και θεωρούμε ένα σύνολο $\{m_1, m_2, \dots, m_k\}$ που παράγει το $M$. Αν το $M$ είναι πρότυπο στρέψης, τότε 
    \[
    a_im_i = 0_M
    \]
    για κάποια μη μηδενικά $a_i \in R$. Επειδή ο $R$ είναι ακέραια περιοχή έπεται ότι $a = a_1a_2\cdots{a_k} \neq 0_R$. Επιπλέον, έχουμε 
    \[
    a(r_1m_1 + \cdots + r_km_k) = r_1a_2\cdots{a_k}(\underbrace{a_1m_1}_{=\, 0_M}) + \cdots + r_1a_1\cdots{a_{k-1}}(\underbrace{a_k m_k}_{=\, 0_M}) = 0_M
    \]
    για κάθε $r_1, \dots, r_k \in R$. Άρα, $a \in \Ann_R(M)$ και γι αυτό $\Ann_R(M) \neq 0$.

    Για το παράδειγμα, έστω $R=\ZZ$, $p$ ένας πρώτος και 
    \[
    M = \bigoplus_{n \ge 1} \ZZ_{p^n}.
    \]
    Το $M$ είναι πρότυπο στρέψης, όμοια με το αντίστοιχο $\ZZ$-πρότυπο του Ερωτήματος (α) (γιατί;). Αν $\Ann_{\ZZ}(M) \neq 0$, τότε θα υπήρχε ένα $k \ge 1$, τέτοιο ώστε 
    \[
    p^km = 0_M
    \]
    \emph{για κάθε} $m \in M$ (γιατί;). Αυτό όμως είναι αδύνατο, διότι 
    \[
    p^k (0 + p\ZZ, \ 0 + p^2\ZZ, \ \dots, \ 0 + p^k\ZZ, \ 1 + p^{k+1}\ZZ,\  0 + p^{k+2}\ZZ, \ \dots) \neq 0_M.
    \]

    \item[(γ)] Υποθέτουμε ότι ο $R$ είναι ακέραια περιοχή και θεωρούμε $\{e_\lambda : \lambda \in \Lambda\}$ μια βάση του $M$. Έστω $m \in M_\tor$ ένα στοιχείο στρέψης του $M$. Τότε 
    \[
    m = \sum_{\lambda \in \Lambda} r_\lambda e_\lambda
    \]
    για κάποια $r_\lambda = 0$ εκτός από ένα το πολύ πεπερασμένο πλήθος και 
    \[
    rm = 0_M
    \]
    για κάποιο μη μηδενικό $r \in R$. Οπότε, 
    \[
    \sum_{\lambda \in \Lambda} rr_\lambda e_\lambda = 0_M
    \]
    και κατά συνέπεια $rr_\lambda = 0_R$, για κάθε $\lambda \in \Lambda$, λόγω της $R$-γραμμικής ανεξαρτησίας των $e_\lambda$. Επειδή ο $R$ είναι ακέραια περιοχή και $r \neq 0_R$, έπεται ότι $r_\lambda = 0_R$, για κάθε $\lambda \in \Lambda$ και γι αυτό $m = 0_M$. Με άλλα λόγια, $M_\tor = 0$ και γι αυτό το $M$ είναι ελεύθερο στρέψης.
  \end{itemize}
\end{proof}

\begin{exercise}{\rm(\cite[Ασκήσεις~10.2.5~και~10.2.6]{DF04})}
  \label{ex:DF_10.2.5-10.2.6}
  Να υπολογιστεί το $\Hom_\ZZ(\ZZ_{30},\ZZ_{21})$ και γενι\-κότερα το $\Hom_{\ZZ}(\ZZ_n, \ZZ_m)$ για κάθε θετικούς ακέραιους $n, m$.
\end{exercise}

\begin{proof}[Λύση]
  Επειδή το $\ZZ_{30}$ είναι κυκλικό $\ZZ$-πρότυπο το οποίο παράγεται από το $1 + 30\ZZ$, κάθε $\phi \in \Hom_{\ZZ}(\ZZ_{30},\ZZ_{21})$ καθορίζεται πλήρως από την τιμή $\phi(1 + 30\ZZ) = a + 21\ZZ$. Επειδή η $\phi$ είναι $\ZZ$-γραμμική έχουμε 
  \begin{align*}
  0 + 21\ZZ &= \phi(0 + 30\ZZ) \\
  &= \phi(30 + 30\ZZ) \\
  &= 30\phi(1 + 30\ZZ) \\
  &= 30(a + 21\ZZ) \\
  &= 30a + 21\ZZ \\
  &= 9a + 21\ZZ.
  \end{align*}
  Με άλλα λόγια, πρέπει να ισχύει ότι 
  \[
  9a \equiv 0 \pmod{21}.
  \]
  Οι λύσεις αυτής της διοφαντικής εξίσωσης δίνονται από τα πολλαπλάσια του $21/\gcd(9,21) = 21/3 = 7$ modulo $21$, δηλαδή $0, 7$ και $14$. Άρα,
  \[
  \Hom_{\ZZ}(\ZZ_{30},\ZZ_{21}) = \{\phi_0, \phi_7, \phi_{14}\}
  \]
  όπου $\phi_n(x + 30\ZZ) = nx + 21\ZZ$. Παρατηρούμε ότι 
  \[
  \Hom_{\ZZ}(\ZZ_{30},\ZZ_{21}) \cong \ZZ_{\gcd(30,21)} = \ZZ_3
  \]
  ως $\ZZ$-πρότυπα και γενικότερα 
  \[
  \Hom_{\ZZ}(\ZZ_n, \ZZ_m) \cong \ZZ_{\gcd(n,m)}
  \]
  ως $\ZZ$-πρότυπα (γιατί;).
\end{proof}

\begin{thebibliography}{99}
    \bibitem{DF04}
    D. S. Dummit και R. M. Foote, 
    \textit{Abstract algebra},
    third edition, Wiley, 2004. 
    % \bibitem{Mal08}
    % Μ. Μαλιάκας, 
    % \textit{Εισαγωγή στην Μεταθετική Άλγεβρα}, 
    % Εκδόσεις Σοφία, 2008.
\end{thebibliography}
\end{document}